% Borrador estructurado para la Tarea 2. Compilar desde 05_T2/ para resolver
% las figuras relativas. Actualizado el 22 de octubre de 2026: incorpora la
% flota con aviso compartido (seccion VI), el ejemplo de base de datos
% regenerado tras corregir la ventana de focos, la auditoria de las fuentes y
% el trabajo pendiente. Falta la corrida del conjunto reservado de evaluacion.
\documentclass[journal,10pt]{IEEEtran}

\usepackage[utf8]{inputenc}
\usepackage[T1]{fontenc}
\usepackage[spanish,es-tabla]{babel}
\usepackage{amsmath,amssymb}
\usepackage{booktabs,tabularx,array}
\usepackage{graphicx}
\usepackage{listings}
\usepackage{xcolor}
\usepackage[hyphens]{url}
\graphicspath{{BD/}}

\lstdefinestyle{codigo}{
  basicstyle=\ttfamily\scriptsize,
  breaklines=true,
  frame=single,
  columns=fullflexible,
  showstringspaces=false,
  backgroundcolor=\color{gray!6},
  belowcaptionskip=6pt
}
\lstset{style=codigo}
\renewcommand{\lstlistingname}{Listado}
% El \@makecaption de IEEEtran no deja espacio bajo el título, y el marco del
% listado, que se dibuja hacia arriba, tapaba su última línea.
\makeatletter
\def\lst@makecaption#1#2{%
  \vbox{\hsize=\linewidth\normalfont\footnotesize\noindent
        #1.\nobreakspace\nobreakspace#2}%
  \vskip\belowcaptionskip}
\makeatother

\begin{document}

\title{Ruteo físico resiliente de UAV para inspección territorial simulada: estudio geoespacial de referencia en Mykolaiv}
\author{Aaron Pozas Oyarce, Diego Pérez Carrasco y Alex Marambio Leyton}
\maketitle

\begin{abstract}
Se implementa y evalúa un mecanismo de ruteo físico resiliente para una misión
\emph{simulada} de inspección de puntos en un recorte urbano de Mykolaiv,
Ucrania. Un único grafo geolocalizado representa lugares de navegación y
segmentos candidatos de vuelo, no enlaces entre UAV. El método genera rutas
mediante Yen, selecciona un par por longitud, riesgo contextual y separación
espacial, y recalcula el trayecto al descubrir un bloqueo simulado. Se
contrasta con rutas basadas en Dijkstra; \texttt{pgr\_dijkstra} y
\texttt{pgr\_KSP} también se usan en el ejemplo de base de datos.
OpenStreetMap, elevación y meteorología histórica alimentan atributos del
entorno; GPSJam, focos térmicos clasificados y eventos UCDP son indicadores
contextuales, no mediciones de tramos intransitables. En 100 escenarios
sintéticos del
conjunto de ajuste con cinco fallas, la propuesta completó el 97\% de las
misiones frente al 87\% de una línea base Dijkstra de par fijo. Dijkstra con
replanificación completó el 96\%; por tanto, con un solo vehículo la ventaja
específica de la selección resiliente no queda demostrada. Se extiende después
el experimento a una \emph{flota reactiva} de tres a cinco UAV que comparten un
aviso de interferencia: el que la detecta lo difunde y el aviso se retransmite
por multisalto entre vehículos, mientras que el que queda dentro de la zona deja
de emitir su \emph{beacon} y su silencio también alerta a la flota; quien recibe
el aviso replanifica en pleno vuelo. Este mecanismo de alerta recoge el efecto
de CLUN-LSR sobre el ruteo físico, sin simular el protocolo. Cuando la
interferencia no siempre se nota antes de entrar en la zona
($p_{\rm det}\leq 0{,}5$), el aviso eleva la tasa de finalización de los UAV
seguidores entre 10 y 17 puntos porcentuales y ningún vehículo empeora por
recibirlo; con detección perfecta en el borde el aviso resulta casi irrelevante
(+0,2 a +0,5 pp). Los resultados proceden del conjunto de ajuste: la corrida
sobre el conjunto reservado sigue pendiente.
\end{abstract}

\begin{IEEEkeywords}
UAV, ruteo resiliente, flota, aviso de interferencia, Yen, Dijkstra, pgRouting,
PostGIS, inspección territorial.
\end{IEEEkeywords}

\IEEEpeerreviewmaketitle

\section{Introducción y delimitación del enfoque}

La primera entrega estudió CLUN-LSR, un protocolo que busca recuperar
\emph{caminos de comunicación} cuando la interferencia afecta los enlaces
de radio de una red de UAV \cite{zheng}; en ese modelo, cada vehículo es un
nodo, cada enlace es una arista y una ``ruta'' indica por cuáles vehículos
pasan los datos. El análisis posterior del planteamiento inicial permitió
identificar que la frecuencia del problema no estaba cuantificada, que la
hipótesis no anticipaba un resultado medible y que los objetivos se
concentraban en preparar datos en lugar de evaluar el mecanismo de ruteo.
También se advirtió que las posiciones de aviones comerciales de OpenSky no
representaban vuelos UAV a baja altura, por lo que esa fuente se descartó
como infraestructura de navegación; hacía falta, además, delimitar una
misión concreta y considerar la dimensión vertical del entorno.

En consecuencia, el presente trabajo estudia otra clase de ruta: el
\emph{desplazamiento físico} de un UAV simulado entre lugares de
observación, representados como nodos unidos por posibles segmentos de
vuelo. El caso de referencia es una inspección territorial civil en
Mykolaiv, Ucrania. El experimento se plantea en dos niveles: primero un UAV
aislado, que solo conoce lo que observa desde el nodo donde se encuentra, y
después una \emph{flota reactiva} de tres a cinco vehículos que recorren la
misma misión escalonados en el tiempo y comparten lo que descubren. De
CLUN-LSR se conserva únicamente el \emph{efecto} de su mecanismo sobre el
ruteo físico ---que la alerta de interferencia llegue a los demás por
retransmisión multisalto y estos replanifiquen---, sin reproducir sus capas,
su gestión de espectro ni MP-OLSR, y sin atribuirle planificación de
trayectorias. Para generar opciones de desplazamiento se emplean los $K$
caminos simples de Yen \cite{yen}, mientras que Luna \emph{et al.}
constituyen un antecedente de replanificación durante misiones de
inspección \cite{luna}. No se reproducen su método de cobertura poligonal ni
su arquitectura de comunicaciones.

\section{Contexto, problema e hipótesis}

\subsection{Caso de estudio y alcance}

Mykolaiv se utiliza como \emph{referencia histórica y cartográfica
provisional}: una evaluación satelital de UNOSAT, basada en imágenes del
21 de julio de 2022, identificó 198 estructuras con daño visible, aunque
el propio producto advierte que su análisis es preliminar y no fue
validado en terreno \cite{unosat}.

Al 19 de octubre de 2026 también existen reportes de ataques recientes
contra la \emph{ciudad}: según la Policía Nacional de Ucrania, un ataque
ruso con drones del 3 de agosto causó una muerte, siete personas heridas y
daños en viviendas \cite{policiaAgo2026}; el 2 de octubre se informó de
otro ataque que hirió a tres residentes y afectó viviendas, un hotel,
locales y depósitos \cite{policiaOctDos2026}, mientras que durante la
noche del 10 de octubre se documentaron daños en un edificio residencial,
edificios administrativos, oficinas y tiendas \cite{policiaOctDiez2026}.
Estos reportes dan cuenta de daños urbanos, pero no registran fallas de
navegación de UAV ni interferencia en las rutas simuladas aquí.

El PNUD, por su parte, documenta una iniciativa de recuperación de Mykolaiv
y la entrega en 2026 de herramientas para recopilar datos espaciales y
aéreos a una comunidad de la región \cite{undpM,undpDatos}. Ese contexto
permite formular \emph{académicamente} una misión civil de observación
territorial, sin afirmar que la misión simulada haya ocurrido o que una
ruta UAV local haya fallado; el grafo tampoco representa corredores aéreos
autorizados ni un plan de vuelo operativo.

La misión fijada en el software parte de $A$, visita en orden los puntos
obligatorios $P_1,P_2,P_3$ y termina en $B$, de modo que se divide en cuatro
tramos origen--destino: $A\!\to\!P_1$, $P_1\!\to\!P_2$,
$P_2\!\to\!P_3$ y $P_3\!\to\!B$; en cada tramo pueden compararse rutas
alternativas entre los mismos extremos. Esta visita de puntos representa
una inspección discreta y no cobertura continua del área, pues el sistema
no modela cámara, campo visual ni porcentaje de superficie observada. Cada
UAV descubre bloqueos solo en los tramos incidentes al nodo que alcanzó; en
la configuración de flota, lo que un vehículo descubre puede llegar a los
demás por el aviso compartido de la Sección~\ref{sec:flota}, pero ninguno
observa el estado global del grafo.

\subsection{Pregunta, hipótesis y objetivos}

La pregunta es si preparar rutas por corredores físicamente separados y
recalcular el trayecto al descubrir un segmento indisponible permite
completar más misiones que preparar dos rutas únicamente según su
distancia. Como dos caminos pueden utilizar aristas distintas y, aun así,
pasar lo bastante cerca para ser afectados por una misma falla espacial,
se consideran tanto su separación como su longitud; el presupuesto y los
escenarios de falla permanecen iguales para todos los métodos comparados.

La segunda pregunta surge de la reunión con el profesor: con un solo dron,
avisar al resto ``es un problema que no existe''. Si varios UAV recorren la
misma misión de forma escalonada, el primero que entra en una zona de
interferencia se convierte en una señal para los que vienen detrás, ya sea
porque la detecta y la comunica o porque deja de responder. Cabe entonces
preguntarse cuánto aporta ese aviso a la finalización de la misión de los
vehículos siguientes, y si aporta algo \emph{por encima} de que cada uno
replanifique con lo que ve por su cuenta.

\textbf{Hipótesis principal refutable:} con parámetros congelados antes de la
evaluación reservada, la estrategia propuesta $C$ tendrá mayor tasa de
finalización de misión que el par Dijkstra fijo $A$, con diferencia pareada
estadísticamente detectable y sin superar el presupuesto establecido.
Además, $C$ se contrastará con un Dijkstra que también recalcula la ruta
desde la posición actual ($D$), para distinguir el beneficio de
\emph{replanificar} del posible aporte de \emph{elegir rutas separadas y
usar puntajes contextuales}; si $C$ no supera a $D$, la mejora frente a
$A$ no podrá atribuirse a esa selección.

\textbf{Hipótesis de la flota:} en escenarios donde la interferencia no
siempre se advierte antes de entrar en la zona, compartir el aviso aumentará
la tasa de finalización de los UAV \emph{seguidores} respecto de la misma
estrategia sin aviso, y ese aumento crecerá con el número de vehículos. La
hipótesis es refutable en dos sentidos: si la detección en el borde es
perfecta, el efecto debería desaparecer, porque cada UAV esquiva la zona por
su cuenta; y si el aviso llegara tarde o indujera desvíos innecesarios,
algunos seguidores completarían la misión sin aviso y la perderían con él.

\textbf{Objetivo general:} diseñar y evaluar sobre un mismo grafo un mecanismo
de ruteo físico que mantenga la capacidad de completar la inspección de puntos
cuando aparecen fallas espaciales simuladas, tanto con un UAV como con una
flota que comparte avisos de interferencia.

\textbf{Objetivos específicos:} (i) construir y documentar el grafo y sus
atributos geoespaciales; (ii) transformar metadata y registros contextuales
en restricciones y puntajes explícitos, preservando su incertidumbre;
(iii) generar y comparar al menos dos rutas por método; (iv) medir la
finalización de misión, la disponibilidad del par, el recorrido y el tiempo
de cálculo en escenarios pareados; y (v) implementar la difusión del aviso
entre UAV y la replanificación en vuelo, y cuantificar su aporte aislándolo
de la replanificación individual.

\section{Infraestructura, ETL y base de datos}

El \emph{espacio de navegación} se representa mediante un único grafo
$G=(V,E)$, donde $V$ contiene lugares geolocalizados por los que podría
pasar el UAV y $E$ contiene los segmentos candidatos de desplazamiento;
estas aristas no son enlaces de radio entre vehículos. Decir que el UAV
está en un nodo significa que alcanzó ese lugar, mientras que su posición
y la distancia ya recorrida pertenecen al estado de la simulación. Aunque
los puntos y segmentos del grafo permanecen definidos, una arista puede
dejar de ser \emph{admisible} cuando aparece una restricción conocida o se
descubre una falla, lo que representa el cambio durante la misión sin
tener que construir un grafo nuevo en cada instante.

La malla reproducible del recorte Mykolaiv contiene 196 nodos y 618
aristas; parte de puntos separados nominalmente por 0,5 km, que pueden
conectarse con hasta ocho vecinos. Sus coordenadas se perturban en una
fracción de 0,25 del tamaño de celda y se retira el 12\% de las aristas
sin desconectar el conjunto (semilla 7), con lo cual se generan
alternativas comparables en el mismo espacio. La existencia de cada
segmento sigue siendo un \emph{supuesto experimental} y no demuestra que
allí exista un corredor aéreo permitido o transitable por un UAV real.

La base utiliza PostgreSQL para almacenar los datos, PostGIS para consultar
su posición espacial y pgRouting para calcular caminos: \texttt{nodos}
guarda coordenadas y elevación; \texttt{aristas}, extremos, longitud y
atributos de cada segmento; \texttt{mision}, los puntos obligatorios;
\texttt{meta\_*} y \texttt{amz\_*}, los datos externos transformados;
\texttt{fuentes}, su procedencia; y \texttt{rutas}/\texttt{evaluaciones},
las salidas. El esquema y el diagrama están en \texttt{BD/schema.sql} y
\texttt{BD/diagrama\_bd.png}. Como las coordenadas horizontales usan
EPSG:4326 y la elevación del terreno y de los obstáculos se incorpora
como atributo, el modelo es \emph{2,5D}: comprueba el margen vertical de
forma aproximada, sin buscar trayectorias con ascensos, descensos o
maniobras tridimensionales.

Todo se ejecuta con \texttt{docker compose up}, de modo que Python solo
corre dentro de los contenedores: el servicio \texttt{db} aloja la base y
el servicio \texttt{main} ejecuta el archivo raíz \texttt{main.py}, que llama
a \texttt{pipeline.py}. Se crea el esquema, se carga el grafo generado en
JSON, se descargan y transforman las fuentes externas, se asignan atributos
y puntajes a las aristas y se calculan rutas de ejemplo y el diagrama; al
terminar, el mismo \texttt{main.py} levanta el sitio Leaflet. Cada JSON de entrada cuenta con un documento
\texttt{ESTRUCTURA\_*.md} que explica sus campos y con un cargador
asociado; la arquitectura y las consultas se detallan en
\texttt{BD/MODELO\_BD.md}. Esta cadena prepara el experimento, cuyo objeto
de evaluación es la decisión de ruta y la reacción ante fallas, no la
mera carga de registros.

La Tabla~\ref{tab:fuentes} distingue entre \emph{metadata}, que describe
condiciones capaces de modificar la admisibilidad o el costo de un tramo
---por ejemplo, una altura excesiva puede excluirlo y el viento puede
encarecerlo---, y \emph{fuentes de amenazas}, que aportan contexto espacial
histórico. La cercanía de un registro no demuestra que el tramo esté
cerrado o interferido: ese cierre ocurre únicamente en los escenarios de
falla sintética definidos más adelante.

\begin{table*}[!t]
\centering
\caption{Fuentes implementadas y traducción al modelo. ``Indicador'' no equivale a falla observada de un UAV.}
\label{tab:fuentes}
\footnotesize
\begin{tabularx}{\textwidth}{@{}>{\raggedright\arraybackslash}p{2.3cm}>{\raggedright\arraybackslash}p{3.2cm}>{\raggedright\arraybackslash}X>{\raggedright\arraybackslash}p{3.2cm}@{}}
\toprule
\textbf{Fuente} & \textbf{Dato extraído} & \textbf{Uso implementado} & \textbf{Límite principal} \\
\midrule
OSM/Overpass \cite{osm} & Huellas y etiquetas \texttt{height} o \texttt{building:levels} & Altura estimada de edificios próximos al tramo; exclusión si el obstáculo alcanza el margen vertical. & Alturas incompletas; tres metros por nivel es un supuesto. \\
Open-Meteo Elevation \cite{elev} & Elevación de nodos y puntos medios, Copernicus GLO-90 & Relieve relativo al punto de salida en la regla 2,5D. & Resolución de 90 m y muestreo de tres puntos por arista. \\
Open-Meteo Historical Weather \cite{meteo} & Serie horaria de viento en un punto & Límite global de viento y factor direccional de costo. & Resolución regional; el script no fija explícitamente el modelo meteorológico en la petición. \\
GPSJam \cite{gpsjam} & Conteos ADS-B de precisión GNSS por celda H3 & Fracción histórica de días con celda en rojo, usada como indicador regional. & Aviones en altura; no detecta jamming UAV ni tiene detalle por calle. \\
The Economist/NASA FIRMS \cite{warfire} & Focos térmicos clasificados como relacionados con la guerra & Densidad histórica para el puntaje; eventos posteriores para centros de fallas sintéticas. & La clasificación es incierta; un foco no bloquea automáticamente una arista. \\
UCDP GED \cite{ucdp} & Eventos georreferenciados y precisión del lugar & Densidad espacial histórica para un puntaje contextual. & \texttt{where\_prec}=2 admite aprox. 25 km de incertidumbre; no mide una barrera de vuelo. \\
\bottomrule
\end{tabularx}
\end{table*}

En la tabla, \emph{área de estudio} designa el recorte urbano donde se
construyó el grafo. OSM (OpenStreetMap) aporta geometrías y etiquetas de
edificios, aunque no entrega alturas verificadas para todos ellos;
GPSJam resume en celdas H3 la precisión de posiciones ADS-B reportadas por
aviones y, por tanto, funciona solo como señal regional e indirecta de
posibles problemas GNSS. UCDP GED reúne eventos de conflicto cuya
localización puede ser imprecisa, mientras que los focos térmicos
clasificados a partir de NASA FIRMS sirven como antecedentes espaciales
cuya relación con un bloqueo de vuelo no está demostrada. Ninguna de
estas fuentes constituye un sensor embarcado en el UAV simulado.

\section{Modelo de restricciones y fallas}
\label{sec:fallas}

El UAV simulado conserva una altura fija respecto de la elevación del punto
de partida $A$; para estimar el margen vertical disponible en una arista
$e$, se compara esa altura de vuelo con la mayor elevación u obstáculo
estimado a lo largo del segmento. La altura relativa del obstáculo se
aproxima mediante
\begin{equation}
 H_e=(z_{\max,e}-z_A)+h_{\max,e},
\end{equation}
donde $z_{\max,e}$ es el máximo de tres muestras de elevación del terreno,
$z_A$ es la elevación del despegue y $h_{\max,e}$ es la mayor altura
estimada de un edificio próximo. La arista $e$ se excluye cuando
$H_e\geq h_{\rm vuelo}-m$, donde $m$ evita considerar suficiente un paso
justo sobre el obstáculo; si falta algún dato, el tramo permanece abierto
en el modelo, sin que eso certifique su seguridad. Además, el modelo
digital de elevación (DEM) puede incorporar parte de la superficie
construida, de modo que sumarle la altura del edificio podría sobreestimar
el obstáculo.

El viento histórico representa una condición general del área y no una
medición instantánea a lo largo de toda la ruta: si supera el umbral
configurado, la misión se marca como no admisible; en caso contrario, su
dirección modifica el costo relativo de cada segmento. Conviene señalar que
la comparación se hace contra la \emph{mediana} de la ventana, 6,43 m/s a la
altura de vuelo frente a un umbral de 12 m/s, de modo que esta restricción
no llega a activarse con los datos disponibles; la ráfaga máxima registrada,
19,1 m/s a 10 m, sí lo superaría. Como el sitio aún puede mostrar rutas de
referencia cuando se excede el umbral, esa pantalla no debe interpretarse
como autorización de vuelo. La interfaz también ofrece una opción de zonas
restringidas, pero no se cargó una fuente que active tal atributo ---el
campo correspondiente es falso en los 618 segmentos--- y, por ello, no se
considera una restricción evaluada.

Cada escenario contiene $n_f$ zonas circulares descritas por
$f=(x_f,y_f,r_f,s_f)$: el centro $(x_f,y_f)$ y el radio $r_f$ delimitan
la zona afectada, mientras que $s_f$ indica la distancia acumulada por el
UAV cuando se activa la falla. Desde ese momento, un segmento que
intersecte el disco queda bloqueado; su centro puede elegirse cerca de un
foco contextual o en otra parte del área, pero el disco \emph{no}
representa un evento observado. Conviene precisar hasta dónde llega ese
vínculo con los focos en la configuración actual: en el experimento de
\texttt{resultados/} los focos son sintéticos y se sortean dentro de la
malla, mientras que en la ruta que usa la base de datos los quince
incendios del período de evaluación caen \emph{todos} fuera de la caja de
la malla, entre 1,95 y 5,12 km de la red. Como el muestreo recorta la
posición a esa caja, el centro de toda falla de origen ``foco'' termina
pegado al borde. Dicho de otro modo, en esa ruta la siembra cerca de un
incendio no se materializa, y así se informa en la documentación técnica
como defecto pendiente. Tanto el radio como la activación y el
bloqueo son supuestos que permiten ensayar cambios de disponibilidad
posteriores a la planificación sin atribuirles una causa real.

El UAV solo descubre el estado de los segmentos conectados al nodo que
acaba de alcanzar, por lo que una falla lejana puede existir sin que la
conozca todavía; al detectarla, actualiza su información y decide si
replanifica. Los discos que cubren un punto obligatorio se vuelven a
muestrear, pues impedirían completar la misión a todos los métodos, de
modo que la tasa de éxito queda \emph{condicionada} por ese criterio.
Ninguna fuente contextual prueba por sí sola el cierre real de una ruta,
ni la pérdida de un reporte permite identificar jamming como su causa.

\begin{table}[!t]
\centering
\caption{Parámetros principales fijados en el código.}
\label{tab:parametros}
\footnotesize
\begin{tabular}{@{}p{3.7cm}p{3.2cm}@{}}
\toprule
\textbf{Parámetro} & \textbf{Valor por defecto} \\
\midrule
Altura de vuelo; margen vertical & 60 m; 10 m \\
Viento máximo; crucero & 12 m/s; 12 m/s \\
Desvío máximo $\alpha$; presupuesto extra $\beta$ & 0,35; 0,60 \\
Pesos $w_L,w_S,w_R$; candidatas Yen & 1, 1, 1; $K=8$ \\
Separación espacial $d_{\rm sep}$ & 0,4 km \\
Radio de incertidumbre $\rho$; penalización $\gamma$ & 0,6 km; 3 \\
Fallas por escenario; radios & 3; 0,25--0,60 km \\
Probabilidad de foco; dispersión & 0,5; 0,5 km \\
\midrule
\multicolumn{2}{@{}l}{\emph{Flota (Sección~\ref{sec:flota})}} \\
UAV por flota $N$; separación de despegues & 4; 1 km \\
Alcance de radio UAV--UAV & 3 km \\
Beacon; beacons para declarar silencio & 2 s; 3 \\
Probabilidad de detección $p_{\rm det}$ & 0,5 \\
\bottomrule
\end{tabular}
\end{table}

\section{Algoritmos de rutas}

Los parámetros de la Tabla~\ref{tab:parametros} controlan decisiones
distintas: $K$ limita las opciones generadas por Yen; $\alpha$ acota el
alargamiento de cada candidata respecto del camino mínimo; $\beta$ fija
los kilómetros adicionales disponibles para toda la misión; y
$d_{\rm sep}$ determina cuándo dos rutas atraviesan un mismo corredor.
Ante un tramo bloqueado, $\rho$ delimita el entorno considerado incierto
y $\gamma$ eleva temporalmente el costo de los tramos cercanos. Todos
estos valores son elecciones de simulación y no especificaciones
certificadas de un UAV real.

\subsection{Línea base y método propuesto}

La \emph{línea base} busca con \texttt{pgr\_dijkstra} el camino de menor
distancia, usando como costo solo los kilómetros de cada arista, y ejecuta
una segunda búsqueda sin las aristas del primer camino para obtener otra
opción. Esto permite compararla con la propuesta sin incorporar puntajes
de amenaza ni separación espacial; sin embargo, si la exclusión
desconecta origen y destino, el código recurre a multiplicar por diez el
costo de las aristas iniciales. \emph{En ese caso el costo deja de ser
exclusivamente kilométrico}, por lo que debe informarse por separado.
El ejemplo obtuvo dos rutas, aunque ello no garantiza dos alternativas
independientes para cualquier par de nodos.

El método propuesto usa Yen para enumerar hasta ocho caminos simples, es
decir, rutas que no repiten nodos \cite{yen,pgr}; como las dos primeras
opciones podrían atravesar el mismo sector, Yen por sí solo no garantiza
resiliencia. Por eso se incorpora la ruta mínima en kilómetros si no
aparece entre las candidatas y se prefieren aquellas que respetan el
desvío y el presupuesto disponible; cuando ninguna cabe, se conserva la
mínima para intentar avanzar, aunque la misión pueda fracasar por exceso
de recorrido. Entre los pares restantes $(p,q)$ se minimiza
\begin{equation}
\begin{aligned}
 J(p,q)={}&w_L\!\left(\frac{L(p)+L(q)}{2L_{\min}}-1\right)\\
          &+w_S S(p,q)+w_R\frac{R(p)+R(q)}{2},
\end{aligned}
\end{equation}
donde $L(p)$ y $L(q)$ son las longitudes de las rutas, $L_{\min}$ la longitud
del camino más corto, $S(p,q)$ el solape espacial y $R(p)$ y $R(q)$ sus
puntajes medios de riesgo contextual. Los tres términos penalizan,
respectivamente, el desvío, el paso de ambas rutas por lugares próximos
aunque no compartan aristas, y la presencia de más indicadores
históricos; así se favorece una alternativa físicamente separada sin
admitir recorridos arbitrariamente largos. Por su parte, el costo para
\emph{buscar} cada candidata combina distancia, viento y riesgo:
\begin{equation}
 c(e,u\!\to\!v)=L_e f_{\rm viento}(u,v)(1+w_R r_e)\pi_e.
\end{equation}
En esta expresión, $L_e$ es la longitud real del tramo,
$f_{\rm viento}$ aumenta el costo si el viento es desfavorable, $r_e$
resume el contexto y $\pi_e$ penaliza el entorno de un bloqueo
descubierto; por tanto, el resultado guía la \emph{elección de ruta} y no
equivale a los kilómetros recorridos. En la base de datos, los
componentes de $r_e$ se agregan mediante $1-\prod_i(1-r_i)$, que produce
un puntaje entre cero y uno sin convertir los registros en una
probabilidad calibrada de fracaso, pues las fuentes son heterogéneas y
no se ha demostrado su independencia.

\begin{lstlisting}[caption={Pseudocódigo del mecanismo implementado; cada tramo une dos puntos consecutivos de la misión.},label={alg:propuesta}]
PLANIFICAR(s, t, conocidas, km_restantes):
  G_adm <- G sin restricciones ni fallas conocidas
  r_min <- Dijkstra(G_adm, s, t, costo = km)
  C <- Yen(G_adm, s, t, K=8, costo = c)
  C <- C union {r_min}
  C <- candidatas que cumplen desvio y holgura
  (p,q) <- par de C con menor J(p,q)
  devolver principal = menor costo c; alternativa = otra

AL_DESCUBRIR(arista bloqueada b):
  registrar b como conocida
  penalizar aristas cercanas a b dentro de rho

REACCIONAR(posicion, siguiente punto t):
  r <- Dijkstra(posicion, t, costo = c,
                sin aristas conocidas bloqueadas)
  si largo(r) excede holgura:
      r <- Dijkstra(posicion, t, costo = km,
                    sin aristas conocidas bloqueadas)
  si no existe r: registrar fracaso
  devolver r
\end{lstlisting}

Yen genera las opciones, mientras que la selección del par y la reacción
ante un bloqueo pertenecen al mecanismo implementado: cuando se descubre
una arista cerrada, el simulador la excluye y encarece los tramos cercanos
por la incertidumbre sobre la extensión de la falla; si el desvío ya no
cabe en el presupuesto, busca el camino disponible más corto. La
alternativa precalculada permite medir la calidad del par, pero no es
necesariamente la que seguirá el UAV, pues $C$ calcula un nuevo Dijkstra
tras el bloqueo; por ello, una mejora frente a un par fijo podría deberse
a la replanificación y no a la elección inicial.

\section{Flota reactiva con aviso compartido}
\label{sec:flota}

Hasta aquí cada UAV solo dispone de lo que observa, de modo que la
información sobre una zona de interferencia muere con el vehículo que la
encuentra. La extensión implementada en \texttt{ruteo/flota.py} hace que esa
información circule: $N\in\{3,4,5\}$ UAV recorren la misma misión
$A\!\to\!P_1\!\to\!P_2\!\to\!P_3\!\to\!B$, despegan de $A$ separados por
1 km recorrido a velocidad de crucero ---unos 83 s a 12 m/s--- y cada uno
conserva su propio presupuesto $(1+\beta)L^*$. El reloj de la simulación se
mide en kilómetros a esa velocidad, por lo que con $N=1$ el simulador de
flota reproduce exactamente la misión individual descrita antes; esa
equivalencia se comprueba en \texttt{tests/test\_flota.py}.

\subsection{Detección y las dos vías de aviso}

Al alcanzar un nodo, el UAV revisa los tramos incidentes y advierte con
probabilidad $p_{\rm det}$ aquellos que cruzan una zona activa. Si no la
advierte y su plan usa uno de esos tramos, entra en la zona y pierde el
enlace: ese vehículo se da por perdido y no se recupera. Sobre esa base
operan dos mecanismos complementarios.

El \emph{aviso por detección} corresponde al efecto que se recoge de
CLUN-LSR. El UAV que advierte la interferencia difunde de inmediato los
tramos afectados, siempre que conserve el enlace, es decir, siempre que no
esté rodeado por la zona. El aviso se retransmite de vehículo en vehículo
---y también a través de la estación en $A$, que actúa como relé--- entre
quienes se encuentran a no más de 3 km entre sí, de manera que un UAV fuera
de alcance lo recibe al reconectarse. No se modelan pérdida de paquetes,
latencia ni ocupación del canal: el alcance se representa como un disco
ideal y la difusión se resuelve calculando las componentes conexas del
grafo de proximidad en cada instante.

El \emph{aviso por silencio} cubre precisamente el caso que el anterior no
puede cubrir: el vehículo que queda dentro de la zona no puede avisar. Cada
UAV emite un \emph{beacon} cada 2 s; cuando deja de emitirlo, quienes
estaban en alcance de su última posición declaran su silencio tras tres
beacons perdidos ---6 s--- y marcan como zona a evitar los tramos situados
a no más de $\rho$ de esa posición. Si al excluirlos no queda ninguna ruta
hacia el punto siguiente, la marca se ignora, de modo que una sospecha
nunca deja a un UAV sin plan. Así, la pérdida de reporte funciona como
indicio de interferencia sin que el modelo afirme que esa es su causa.

Quien recibe un aviso incorpora los tramos afectados a su información:
$D_A$ los excluye y $C_A$, además, penaliza su entorno mediante $\pi_e$. Si
la parte pendiente de su ruta los atraviesa, \textbf{replanifica en ese
instante} desde el nodo hacia el que vuela, sin esperar a toparse con el
bloqueo; un UAV que aún no despega planifica ya con el aviso incorporado.
No se admite girar en mitad de un tramo, de modo que el cambio rige desde
el nodo siguiente.

\begin{lstlisting}[caption={Difusión del aviso y reacción en vuelo; complementa el pseudocódigo anterior.},label={alg:flota}]
AL_DETECTAR(uav, tramos bloqueados B, t):
  registrar B en el conocimiento de uav
  si uav NO esta rodeado por la zona:
      emitir aviso(deteccion, B, t)

AL_PERDER(uav, nodo, t):
  programar silencio(uav, nodo) en t + 3*beacon

SILENCIO(uav, nodo, t):
  oyentes <- UAV conectados con nodo en t
  si oyentes no vacio:
      emitir aviso(silencio, evitar = tramos a <= rho de nodo, t)

DIFUNDIR(t):
  para cada componente conexa por alcance de radio:
      todos sus miembros conocen la union de sus avisos

AL_RECIBIR(uav, aviso, t):
  actualizar conocidas y evitar
  si la ruta pendiente de uav cruza lo avisado:
      replanificar desde el nodo hacia el que vuela
\end{lstlisting}

\subsection{Estrategias comparadas y métricas}

Para separar el aporte del aviso del aporte de la selección de rutas se
comparan cinco configuraciones sobre los \emph{mismos} escenarios: $A$, $D$
y $C$ sin aviso, donde cada UAV aprende solo lo que ve, y $D_A$ y $C_A$ con
aviso compartido. La diferencia $D_A-D$ mide entonces el valor del aviso
sobre una misma política de replanificación, mientras que $C_A-D_A$ mide lo
que añade la selección resiliente por encima del aviso. El sorteo de
detección se fija por escenario, vehículo y tramo, así que todas las
estrategias enfrentan exactamente las mismas oportunidades de advertir una
zona.

La métrica principal es $R_{\rm seguidores}$, la tasa de finalización de los
UAV $2\ldots N$, que son los únicos que pueden aprovechar un aviso ajeno; el
líder se informa por separado, porque su desempeño no debería cambiar con el
mecanismo. Se registran además la tasa sobre todos los vehículos
$R_{\rm uav}$, la fracción de escenarios en que al menos uno completa la
misión, los UAV perdidos por escenario, los cambios de ruta provocados por un
aviso y la demora entre la detección ---o la pérdida del vehículo--- y el
instante en que cada uno de los demás se entera. Como los UAV de un mismo
escenario comparten fallas y no son independientes, el $p$ de McNemar
calculado sobre pares (escenario, UAV) se informa como orientativo.

\section{Metodología de evaluación}

El presupuesto se fija como $L_{\max}=(1+\beta)L^*$, donde $L^*$ suma las
distancias mínimas \emph{sin fallas} de los cuatro tramos; con
$\beta=0{,}60$ se permite recorrer hasta un 60\% más que ese mínimo para
absorber desvíos, aunque ninguna ruta individual recibe por ello una
holgura ilimitada. Una misión se considera exitosa solo si visita
$P_1,P_2,P_3$ en orden, llega a $B$ y no supera $L_{\max}$, por lo que
la métrica principal es la proporción de escenarios que cumplen todas
esas condiciones:
\begin{equation}
 R_{\rm mision}=\frac{\#\{\text{misiones completadas}\}}
                         {\#\{\text{escenarios evaluados}\}}.
\end{equation}
La métrica complementaria $R_{\rm par}$ expresa la proporción de tramos
y escenarios en que al menos una ruta preparada antes de salir permanece
libre de las fallas simuladas; describe la calidad del \emph{par}, no la
finalización de la misión, porque el UAV puede desconocer una falla al
partir o encontrarla más tarde. También se registran kilómetros
recorridos, sobrelongitud respecto de $L^*$, replanificaciones y tiempo
de cálculo. La diferencia $\Delta R=100(R_C-R_A)$ se informa en
\emph{puntos porcentuales} ---de 87\% a 97\% son $+10$ puntos--- y la
prueba exacta de McNemar compara qué método completa cada uno de los
\emph{mismos} escenarios.

Se comparan $A$ (dos Dijkstra fijos), $B$ (dos caminos más cortos de Yen
sin selección), $D$ (Dijkstra que replanifica), $C$ (propuesta), $C0$
(propuesta sin penalización del entorno) y $C1$ (sin control de
presupuesto). Los pares precalculados de $A$ y $B$ muestran qué ocurre
sin replanificación, $D$ aísla el beneficio de recalcular la ruta y $C0$
y $C1$ examinan dos componentes de $C$; para que la comparación sea
válida, todos comparten grafo, restricciones, puntos, escenarios y límite
de recorrido. Las semillas 0--99 se destinan al \emph{ajuste}, donde se
probaron una, tres y cinco fallas con probabilidades de situarlas cerca
de un foco de 0, 0,5 y 0,8; las semillas 10\,000--10\,299 quedan
reservadas para la evaluación final tras congelar parámetros, y el sitio
usa otras semillas para demostración. Los experimentos de flota emplean esas
mismas semillas de ajuste, de modo que las cinco estrategias con y sin aviso
se enfrentan a escenarios idénticos y la comparación sigue siendo pareada; en
ellos se barren $N\in\{3,4,5\}$ y
$p_{\rm det}\in\{1;\,0{,}5;\,0\}$, por tratarse de un supuesto sin respaldo
empírico disponible.

Hay dos ejecuciones que deben distinguirse: los resultados comparativos de
\texttt{resultados/} provienen del simulador Python, con grafo y focos de
riesgo \emph{sintéticos}, y evalúan la misión completa
$A\!\to\!P_1\!\to\!P_2\!\to\!P_3\!\to\!B$; en cambio, el pipeline
PostGIS/pgRouting carga las seis fuentes externas y produce un ejemplo de
rutas entre el primer y el último nodo de \texttt{mision}. Este segundo
resultado comprueba consultas y permite visualizar un par
origen--destino, pero no evalúa por sí mismo las cuatro visitas
obligatorias; por eso, sus cifras de resiliencia se interpretan por
separado de las del experimento Python.

\section{Resultados disponibles y lectura crítica}

El repositorio conserva en \texttt{resultados/} cifras
\emph{preliminares del conjunto de ajuste}; con la malla y misión por
defecto, $L^*=13{,}1637$ km. La Tabla~\ref{tab:resultados} presenta
cuatro configuraciones de 100 escenarios cada una, en las que todos los
métodos enfrentan las mismas fallas simuladas: $n_f$ indica su número y
$\beta$ el recorrido adicional permitido. Ningún valor de esta tabla
proviene de vuelos reales ni del ejemplo de rutas de la base de datos.

\begin{table}[!t]
\centering
\caption{Tasa $R_{\rm mision}$ en el conjunto de ajuste; probabilidad de foco 0,5.}
\label{tab:resultados}
\footnotesize
\begin{tabular}{@{}ccrrrr@{}}
\toprule
$n_f$ & $\beta$ & $A$ & $B$ & $D$ & $C$ \\
\midrule
3 & 0,60 & 0,91 & 0,82 & 0,96 & 0,97 \\
5 & 0,60 & 0,87 & 0,73 & 0,96 & 0,97 \\
5 & 0,20 & 0,75 & 0,72 & 0,94 & 0,90 \\
5 & 0,10 & 0,70 & 0,71 & 0,91 & 0,84 \\
\bottomrule
\end{tabular}
\end{table}

Con cinco fallas y $\beta=0{,}60$, $C$ completó 97 de 100 misiones frente a
87 de $A$: son $+10$ puntos porcentuales ($p=0{,}01294$ en la prueba
pareada). Sin embargo, $D$ completó 96, casi lo mismo que $C$ ($+1$ punto;
$p=1$). Cuando el presupuesto baja a $\beta=0{,}10$, $C$ queda siete puntos
por debajo de $D$. El par de $C$ tiene mayor $R_{\rm par}$ que el de $A$
en la configuración de cinco fallas y $\beta=0{,}60$ (0,94 frente a 0,9075),
lo que indica que sus
dos rutas precalculadas resisten mejor las fallas simuladas. No demuestra
que esa selección cause la mayor finalización de misiones: el resultado
frente a $D$ sugiere que \emph{replanificar} explica gran parte de la
diferencia frente a $A$. En esa corrida local, el cálculo de $C$ tomó en
promedio unos 34 ms por misión; $A$ y $D$, menos de 2 ms.

\subsection{Flota con aviso compartido}

La Tabla~\ref{tab:flota} reúne el barrido de flota sobre el mismo conjunto
de ajuste, con 100 escenarios por celda. Las tasas corresponden a los UAV
seguidores y los $\Delta R$ se expresan en puntos porcentuales sobre esa
población.

\begin{table}[!t]
\centering
\caption{Flota en el conjunto de ajuste: $R_{\rm seguidores}$ y aporte del aviso. Primera fila con tres fallas; el resto, con cinco. Probabilidad de foco 0,5 y $\beta=0{,}60$.}
\label{tab:flota}
\footnotesize
\begin{tabular}{@{}ccrrrrr@{}}
\toprule
$p_{\rm det}$ & $N$ & $D$ & $D_A$ & $C_A$ & $D_A\!-\!D$ & $C_A\!-\!D_A$ \\
\midrule
0,5 & 4 & 0,770 & 0,870 & 0,930 & $+10{,}0$ & $+6{,}0$ \\
\midrule
1   & 3 & 0,930 & 0,935 & 0,960 & $+0{,}5$  & $+2{,}5$ \\
1   & 4 & 0,947 & 0,950 & 0,967 & $+0{,}3$  & $+1{,}7$ \\
1   & 5 & 0,955 & 0,958 & 0,965 & $+0{,}2$  & $+0{,}8$ \\
0,5 & 3 & 0,690 & 0,805 & 0,835 & $+11{,}5$ & $+3{,}0$ \\
0,5 & 4 & 0,667 & 0,807 & 0,857 & $+14{,}0$ & $+5{,}0$ \\
0,5 & 5 & 0,637 & 0,807 & 0,873 & $+17{,}0$ & $+6{,}5$ \\
0   & 3 & 0,605 & 0,740 & 0,730 & $+13{,}5$ & $-1{,}0$ \\
0   & 4 & 0,573 & 0,720 & 0,740 & $+14{,}7$ & $+2{,}0$ \\
0   & 5 & 0,542 & 0,705 & 0,740 & $+16{,}2$ & $+3{,}5$ \\
\bottomrule
\end{tabular}
\end{table}

Cuatro lecturas se desprenden de esos números. Primero, \textbf{el aviso
nunca perjudicó a un seguidor}: en ninguna celda hubo un UAV que completara
la misión sin aviso y la perdiera con él, de modo que las columnas de
escenarios discordantes son unilaterales ---por ejemplo, 68 contra 0 con
$p_{\rm det}=0{,}5$ y cinco vehículos--- y el $p$ de McNemar queda por debajo
de $10^{-4}$ en todas las celdas con $p_{\rm det}\leq0{,}5$. Segundo, el
beneficio aparece solo cuando la interferencia no es evidente antes de
entrar: con $p_{\rm det}=1$ el aviso aporta entre $+0{,}2$ y $+0{,}5$ pp, pues
cada UAV esquiva la zona por su cuenta y lo único que ahorra el aviso son
kilómetros. Tercero, el aporte crece con el tamaño de la flota ---de $+11{,}5$
a $+17{,}0$ pp al pasar de tres a cinco UAV con $p_{\rm det}=0{,}5$---, lo que
es coherente con el mecanismo: hay más vehículos detrás que aprovechan el
mismo aviso. Cuarto, los UAV perdidos por escenario bajan de forma
consistente, por ejemplo de 1,24 con $D$ a 0,63 con $C_A$ en la celda de
cuatro vehículos y $p_{\rm det}=0{,}5$.

La demora del aviso fue de 0 a 1,8 s por detección ---casi todas las
recepciones ocurren en el mismo instante, y los retrasos corresponden a
vehículos que estaban fuera de alcance y lo recibieron al reconectarse--- y
de 6 s por silencio, el tiempo de declarar tres beacons perdidos. La única
excepción es $D_A$ con cinco vehículos y $p_{\rm det}=0{,}5$, donde la media
sube a 6,7 s por el mismo motivo: algún vehículo no estaba en alcance cuando
se declaró el silencio y recibió el aviso después, por retransmisión.

Conviene subrayar que \textbf{el UAV que se pierde no se salva}: el
beneficio es íntegramente para los que vienen detrás, y por eso la métrica
principal excluye al líder. La selección resiliente por encima del aviso
($C_A-D_A$) suma entre $+3$ y $+6{,}5$ pp con $p_{\rm det}=0{,}5$, pero solo
resulta detectable en dos de esas cuatro celdas: con tres fallas ($p=0{,}006$)
y con cinco fallas y cinco vehículos ($p=0{,}003$). Con tres y cuatro vehículos
no lo es ($+3{,}0$ pp, $p=0{,}42$; $+5{,}0$ pp, $p=0{,}058$), y tampoco con
$p_{\rm det}=0$ ni con $p_{\rm det}=1$. Es el primer escenario en que el
aporte propio llega a distinguirse de la replanificación simple, aunque de
forma parcial; además, proviene del conjunto de ajuste y no de la evaluación
reservada, de modo que no basta para dar por confirmada la hipótesis
principal.

La ejecución completa con \texttt{docker compose up} termina con sus 18 pasos
correctos y permite comprobar el esquema, las consultas de pgRouting y las
seis extracciones; la última verificación, con descarga completa desde cero,
se hizo el 22 de octubre de 2026. Su bitácora documenta 33\,630 edificios OSM
(1\,019 con altura estimada), 814 muestras de elevación, 720 horas de viento,
siete celdas GPSJam, 41 focos térmicos clasificados y 65 eventos UCDP. Las
otras cinco fuentes son estáticas o se filtran por ventanas de fechas fijas,
y devolvieron los mismos conteos que la corrida del 19 de octubre;
OpenStreetMap, en cambio, se edita continuamente: aquella descarga había
traído 33\,593 edificios (1\,022 con altura), sin que la diferencia cambie las
exclusiones ni las rutas de ejemplo. Estos números
describen \emph{registros extraídos o cargados}, no fallas sufridas por UAV. Las
modificaciones posteriores del diagrama de base de datos mejoran su
legibilidad; no cambian el grafo de navegación ni los resultados del
simulador.

La primera corrida reveló un error que afectaba al ejemplo de base de datos.
Los focos del período de evaluación estaban en un margen de 5 km alrededor
del área de estudio, pero la consulta de focos buscaba solo dentro del área
estricta; al no encontrarlos, usaba focos \emph{históricos}, que ya habían
servido para calcular el riesgo de los tramos. La ventana de búsqueda se
amplió en \texttt{comun/resiliencia.py} y el ejemplo se regeneró después del
arreglo: entre los nodos 15 y 181, la línea base entrega rutas de 8,462 y
9,166 km y la propuesta, de 8,462 y 8,587 km, con riesgo medio menor en la
segunda ($0{,}4466$ frente a $0{,}5253$). Sobre 100 escenarios de demostración,
$R_{\rm par}$ resulta 0,91 para Dijkstra y 0,90 para la propuesta, mientras
que $R_{\rm mision}$ da 0,91 para $A$ y 1,00 para $D$ y $C$
($\Delta R=9{,}0$ pp, $p=0{,}00391$). Estas cifras describen un par
origen--destino y no la misión con los tres puntos intermedios, por lo que
se interpretan por separado de la Tabla~\ref{tab:resultados}. Sigue sin
existir una corrida conservada del conjunto reservado de evaluación.

\section{Discusión, límites y cumplimiento de la Tarea 2}

El prototipo dispone de la cadena pedida por la rúbrica: Docker, SQL de
creación, diagrama, extracción a JSON, cargadores, pesos, consultas
\texttt{pgr\_dijkstra}/\texttt{pgr\_KSP}, métricas y un sitio Leaflet para
elegir origen, destino y preferencias, con una sección que anima el mismo
escenario de flota con aviso y sin aviso y permite mover, quitar o añadir
los puntos de observación antes de simular; cuando la secuencia editada se
aparta de la misión almacenada, la propia pantalla advierte que esas cifras
no son comparables con las de este informe. La ejecución completa comprueba
que el pipeline termina y que las cifras del ejemplo son posteriores al
ajuste de la consulta; no sustituye la demostración en vivo de cada función
del sitio. \texttt{ENTREGA\_TAREA2.md} indica dónde encontrar la evidencia
asociada a los 60 puntos. El enunciado también exige alojar la versión
evaluada en GitLab UDP; el remoto de esta copia sigue siendo GitHub y esa
sincronización está pendiente.

Las principales limitaciones científicas son: una malla y una misión
sintéticas; conocimiento local perfecto de aristas incidentes; focos y
fallas unidos por reglas supuestas; cobertura incompleta de alturas OSM;
viento regional; y localizaciones de eventos con precisión variable. Un foco
térmico, un registro de conflicto o la baja precisión GNSS de un avión no
equivalen a un tramo UAV cerrado.

Una auditoría de las seis fuentes muestra además que su aporte es muy
desigual, y conviene declararlo antes que presentarlas como equivalentes.
Los edificios de OpenStreetMap hacen el trabajo decisivo: 79 de las 83
exclusiones de segmentos necesitan su altura, y solo 4 se explican por el
relieve. Los eventos UCDP discriminan bien, con 316 valores distintos entre
los 618 segmentos. En el otro extremo, GPSJam \emph{no discrimina en
absoluto}: las celdas H3 de la fuente miden unos 1\,775 km$^2$ y el área de
estudio ocupa unos 44 km$^2$, de modo que cabe entera en una sola celda y los
618 segmentos reciben el mismo valor, 0,3333. Como el riesgo agregado es una
función monótona de los componentes, esa constante no altera el orden de
riesgo de ningún segmento; únicamente eleva el riesgo medio desde
aproximadamente 0,17 hasta 0,4511, por lo que esa cifra no debe leerse como
una medida de amenaza. Los incendios históricos aportan poco, pues solo 7 de
los 26 quedan dentro del radio de corte de 1,5 km, y las alturas de
edificios cubren solo 1\,019 de 33\,630 registros; como un dato ausente no
bloquea, el resto se trata de hecho como si no tuviera altura.

El trabajo pendiente más inmediato, ordenado por lo que aporta, es: cargar
una fuente de zonas restringidas ---por ejemplo, aeródromos y otras áreas de
acceso restringido etiquetadas en OpenStreetMap, obtenidas en la misma
consulta--- para que esa preferencia deje de ser inerte; retirar GPSJam del
puntaje de riesgo e informarlo como una constante regional; evaluar si la
restricción de viento debe comparar un percentil alto o la ráfaga en lugar de
la mediana; y corregir el recorte de los focos al borde de la malla descrito
en la Sección~\ref{sec:fallas}. Además, el sitio muestra un par
origen--destino y calcula su resiliencia para esa consulta; no debe
rotularse como éxito de la misión completa con los tres puntos intermedios.
Tampoco se ha demostrado que el método propuesto supere consistentemente a
$D$, que también replanifica.

La flota añade sus propias limitaciones. El enlace se representa como un
disco de alcance fijo, sin pérdida de paquetes, latencia ni saturación del
canal, de modo que lo evaluado es el \emph{efecto} del aviso sobre el ruteo
físico y no el desempeño de un protocolo: CLUN-LSR sigue siendo el
antecedente del mecanismo, no algo que este trabajo implemente o replique.
La probabilidad de detección es un parámetro sin respaldo empírico ---no hay
datos públicos sobre con cuánta anticipación un UAV advierte la
interferencia---, y como el resultado depende fuertemente de ella se informa
el barrido completo en lugar de una sola cifra. Todos los vehículos
comparten misión, velocidad y presupuesto, despegan escalonados de forma
regular y no giran en mitad de un tramo; tampoco se modela la coordinación
entre ellos más allá de la difusión del aviso. Por último, los UAV de un
mismo escenario comparten fallas, así que el $p$ de McNemar sobre pares
(escenario, UAV) no cumple el supuesto de independencia y debe leerse como
orientativo. El alcance real del experimento es comprobar si una estrategia
reacciona mejor \emph{bajo el modelo de fallas declarado}, no garantizar su
desempeño en vuelos reales.

Queda fuera del trabajo, y se declara en lugar de omitirse, una línea que
sugirió el profesor: considerar la defensa antiaérea como causa de pérdida
de vehículos. No se encontró una fuente pública adecuada para incorporarla
como dato, y el modelo solo representa zonas abstractas de indisponibilidad.
El aviso por silencio ya es compatible con esa situación, porque reacciona a
la desaparición de un UAV sin suponer su causa; lo que falta es caracterizar
la amenaza ---su alcance y si un desvío la evita---, de modo que el radio de
la zona a evitar no sea un supuesto. Ese análisis se deja como trabajo
futuro.

\section{Conclusiones}

La Tarea 2 cuenta con una implementación coherente de ruteo físico y un
experimento reproducible de inspección de puntos. El modelo distingue las
rutas de comunicación de las trayectorias físicas y utiliza un grafo de
navegación único y una hipótesis cuantitativa. Con un solo vehículo, la
propuesta mejora a un Dijkstra de par fijo en los escenarios de ajuste, pero
no establece una ventaja consistente frente a Dijkstra con replanificación.

El aporte que sí se sostiene con claridad es el de la flota: cuando la
interferencia no se advierte siempre antes de entrar en la zona, compartir
el aviso ---por detección o por el silencio del vehículo que quedó
dentro--- y replanificar en vuelo eleva la finalización de los UAV
seguidores entre 10 y 17 puntos porcentuales, el efecto crece con el número
de vehículos y ningún seguidor empeora por recibir el aviso. El resultado
también delimita cuándo el mecanismo es irrelevante: con detección perfecta
en el borde, cada UAV esquiva la zona por su cuenta. Es, además, el primer
contexto en que la selección resiliente llega a distinguirse de la simple
replanificación, aunque solo en algunas configuraciones y con márgenes
menores, de $+3$ a $+6{,}5$ pp.

Estas conclusiones valen para el conjunto de ajuste. La corrida reservada
con parámetros congelados sigue pendiente y es la que permitirá afirmar o
refutar las dos hipótesis. El alcance de los resultados es la simulación
declarada, no la operación de UAV reales.

\begin{thebibliography}{99}
\bibitem{zheng} S. Zheng, X. Zhu, Z. Qin y C. Dong, ``Cross-layer UAV network routing protocol for spectrum denial environments,'' \emph{Ad Hoc Networks}, 2024, doi: 10.1016/j.adhoc.2024.103702.
\bibitem{yen} J. Y. Yen, ``Finding the K Shortest Loopless Paths in a Network,'' \emph{Management Science}, vol. 17, no. 11, pp. 712--716, 1971, doi: 10.1287/mnsc.17.11.712.
\bibitem{luna} M. A. Luna \emph{et al.}, ``A multi-UAV system for coverage path planning applications with in-flight re-planning capabilities,'' \emph{Journal of Field Robotics}, 2024, doi: 10.1002/rob.22342.
\bibitem{unosat} UNOSAT, ``Mykolaiv Damage Assessment Overview Map,'' 14 oct. 2022. [En línea]. Disponible: \url{https://unosat.org/products/3404}.
\bibitem{policiaAgo2026} Policía Nacional de Ucrania, ``Zahybla ta semero postrazhdalykh, z nykh dvoie ditei: u Mykolaievi politseiski dokumentuiut naslidky cherhovoho voiennoho zlochynu RF,'' 4 ago. 2026. [En línea]. Disponible: \url{https://mk.npu.gov.ua/news/zahybla-ta-semero-postrazhdalykh-z-nykh-dvoie-ditei-u-mykolaievi-politseiski-dokumentuiut-naslidky-cherhovoho-voiennoho-zlochynu-rf}.
\bibitem{policiaOctDos2026} Policía Nacional de Ucrania, ``Postrazhdaly troie hromadian, zaznaly poshkodzhen tsyvilni obiekty: u Mykolaievi politseiski dokumentuiut naslidky cherhovykh voiennykh zlochyniv RF,'' 2 oct. 2026. [En línea]. Disponible: \url{https://npu.gov.ua/news/postrazhdaly-troie-hromadian-zaznaly-poshkodzhen-tsyvilni-obiekty-u-mykolaievi-politseiski-dokumentuiut-naslidky-cherhovykh-voiennykh-zlochyniv-rf}.
\bibitem{policiaOctDiez2026} Policía Nacional de Ucrania, ``Politseiski fiksuiut naslidky nichnoi cherhovoi vorozhoi ataky RF na Mykolaiv,'' 10 oct. 2026. [En línea]. Disponible: \url{https://npu.gov.ua/news/politseiski-fiksuiut-naslidky-nichnoi-cherhovoi-vorozhoi-ataky-rf-na-mykolaiv}.
\bibitem{undpM} PNUD, ``Mykolaiv Strategic Partnership,'' 2025--2028. [En línea]. Disponible: \url{https://www.undp.org/ukraine/projects/mykolaiv-strategic-partnership}.
\bibitem{undpDatos} PNUD, ``UNDP provides equipment to five Ukrainian communities to support data-driven recovery,'' 31 mar. 2026. [En línea]. Disponible: \url{https://www.undp.org/ukraine/press-releases/undp-provides-equipment-five-ukrainian-communities-support-data-driven-recovery}.
\bibitem{osm} OpenStreetMap Contributors, ``Copyright and License.'' [En línea]. Disponible: \url{https://www.openstreetmap.org/copyright}.
\bibitem{elev} Open-Meteo, ``Elevation API: Copernicus DEM GLO-90.'' [En línea]. Disponible: \url{https://open-meteo.com/en/docs/elevation-api}.
\bibitem{meteo} Open-Meteo, ``Historical Weather API.'' [En línea]. Disponible: \url{https://open-meteo.com/en/docs/historical-weather-api}.
\bibitem{gpsjam} GPSJam, ``FAQ: what the map shows and its limitations.'' [En línea]. Disponible: \url{https://gpsjam.org/faq}.
\bibitem{warfire} The Economist, ``The Economist's Ukraine war-fire model: source code and data.'' [En línea]. Disponible: \url{https://github.com/TheEconomist/the-economist-war-fire-model}.
\bibitem{ucdp} Uppsala Conflict Data Program, ``UCDP Georeferenced Event Dataset Codebook.'' [En línea]. Disponible: \url{https://ucdp.uu.se/downloads/ged/ged251.pdf}.
\bibitem{pgr} pgRouting Project, ``pgRouting Manual, version 3.8.'' [En línea]. Disponible: \url{https://docs.pgrouting.org/latest/en/}.
\end{thebibliography}

\end{document}
